% p374_2; source PDF page 374; manually verified editable table.
\begin{tikzpicture}[x=1cm,y=-1cm,inner sep=0pt,outer sep=0pt]
\path[use as bounding box] (0,0) rectangle (16.000,10.210);
\draw[line width=.4pt] (0.000,0.000) rectangle (2.100,0.865);
\node[align=center,text width=1.840cm,font={\linespread{1}\fontsize{10}{13.50}\selectfont\bfseries}] at (1.050,0.432) {虚假诉讼罪};
\draw[line width=.4pt] (2.100,0.000) rectangle (6.000,0.865);
\node[align=center,text width=3.640cm,font={\linespread{1}\fontsize{10}{13.50}\selectfont\bfseries}] at (4.050,0.432) {内容};
\draw[line width=.4pt] (6.000,0.000) rectangle (16.000,0.865);
\node[align=center,text width=9.740cm,font={\linespread{1}\fontsize{10}{13.50}\selectfont\bfseries}] at (11.000,0.432) {具体表现};
\draw[line width=.4pt] (0.000,0.865) rectangle (2.100,3.980);
\node[align=center,text width=1.840cm,font={\linespread{1}\fontsize{10}{13.50}\selectfont}] at (1.050,2.423) {行为方式（成立条件）};
\draw[line width=.4pt] (2.100,0.865) rectangle (6.000,3.980);
\node[align=left,text width=3.640cm,font={\linespread{1}\fontsize{10}{13.50}\selectfont}] at (4.050,2.423) {以捏造的事实提起虚假的民事诉讼。\\{}1.执行程序属于诉讼的一部分；\\{}2.仲裁不属于诉讼};
\draw[line width=.4pt] (6.000,0.865) rectangle (16.000,3.980);
\node[align=left,text width=9.740cm,font={\linespread{1}\fontsize{10}{13.50}\selectfont}] at (11.000,2.423) {1.乙欠甲30万元，甲将欠条修改为“欠80万元”，提起诉讼。多数说认为，“虚假”要求完全假，甲不构成虚假诉讼罪。少数说认为，“虚假”只要求部分假，甲构成虚假诉讼罪。\\{}2.隐瞒债务已经全部清偿的事实，向法院提起诉讼，要求还债。这种隐瞒事实，视为捏造事实，构成虚假诉讼罪};
\draw[line width=.4pt] (0.000,3.980) rectangle (2.100,6.495);
\node[align=center,text width=1.840cm,font={\linespread{1}\fontsize{10}{13.50}\selectfont}] at (1.050,5.238) {行为后果（既遂条件）};
\draw[line width=.4pt] (2.100,3.980) rectangle (6.000,6.495);
\node[align=left,text width=3.640cm,font={\linespread{1}\fontsize{10}{13.50}\selectfont}] at (4.050,5.238) {妨害司法秩序或者严重侵害他人合法权益};
\draw[line width=.4pt] (6.000,3.980) rectangle (16.000,6.495);
\node[align=left,text width=9.740cm,font={\linespread{1}\fontsize{10}{13.50}\selectfont}] at (11.000,5.238) {1.致使法院基于捏造的事实\CJKunderwave{采取财产保全或者行为保全措施}的；\\{}2.\CJKunderwave{致使法院开庭审理}。\\{}满足任一条件，便构成既遂。（2023年试题）\\{}（少数说：法院受理便既遂）};
\draw[line width=.4pt] (0.000,6.495) rectangle (2.100,10.210);
\node[align=center,text width=1.840cm,font={\linespread{1}\fontsize{10}{13.50}\selectfont}] at (1.050,8.353) {罪数};
\draw[line width=.4pt] (2.100,6.495) rectangle (16.000,10.210);
\node[align=left,text width=13.640cm,font={\linespread{1}\fontsize{10}{13.50}\selectfont}] at (9.050,8.353) {1.甲以非法占有乙的财物为目的，提起虚假民事诉讼，法官受骗，判决甲胜诉，乙遭受财产损失。甲的一个行为同时触犯虚假诉讼罪和诈骗罪，想象竞合，择一重罪论处。\\{}2.甲与法官勾结，以非法占有乙的财物为目的，提起虚假民事诉讼，法官故意判决甲胜诉，乙遭受财产损失。甲与法官触犯三个罪的共同犯罪：针对司法秩序法益，触犯虚假诉讼罪、民事枉法裁判罪；针对被害人乙的财产，触犯盗窃罪（因为非法转移占有乙的财物）。想象竞合，择一重罪论处（2021年试题）};
\end{tikzpicture}
